% Copyright (C) 2026 Christopher Hoover (ch@murgatroid.com)
% SPDX-License-Identifier: CC-BY-NC-SA-4.0
% Licensed under Creative Commons Attribution-NonCommercial-ShareAlike 4.0 International
% https://creativecommons.org/licenses/by-nc-sa/4.0/
\documentclass[12pt]{article}
\usepackage[T1]{fontenc}
\usepackage{lmodern}
\usepackage{amsmath}
\usepackage{geometry}
\usepackage{siunitx}
\usepackage{microtype}
\usepackage[hidelinks]{hyperref}
\usepackage{bookmark}

\geometry{margin=1in}

\hypersetup{
    pdftitle={A Diode Is Not a Square-Law Device},
    pdfauthor={Christopher Hoover (AI6KG, ch@murgatroid.com)},
    pdfsubject={Diode square-law approximation, detection, mixing}
}

\title{A Diode Is Not a Square-Law Device}
\author{Christopher Hoover, AI6KG \\ \texttt{ch@murgatroid.com}\thanks{\copyright{} 2026 Christopher Hoover. Licensed under CC BY-NC-SA 4.0: \url{https://creativecommons.org/licenses/by-nc-sa/4.0/}}}
\date{}

\begin{document}

\maketitle

Pick up almost any RF textbook and you will find the diode described as a ``square-law device.''
Crystal radio enthusiasts use it to explain why a detector converts RF to audio.
Mixer designers use it to explain frequency conversion.
It appears in receiver noise analysis, intermodulation calculations, and detector sensitivity specifications.

There is just one problem: it is not true.

A diode's current--voltage characteristic is exponential, not quadratic.
Understanding why we use the square-law approximation anyway---and when it breaks down---matters in practice.

\section*{What the diode actually does}

The current through a diode rises exponentially with forward voltage:
\begin{equation*}
  I \approx I_S \, e^{V / n V_T},
\end{equation*}
where $I_S$ is a small leakage current, $n$ is a device-dependent ideality factor near 1--2,
and $V_T \approx \SI{26}{\milli\volt}$ at room temperature is the thermal voltage.
On a linear scale the I--V curve shoots nearly straight up once conduction begins;
on a logarithmic current scale it is a straight line spanning many decades.
It is not a parabola.

\section*{Why ``square law'' seems right}

Apply a small AC signal $v$ on top of a DC bias and the diode's AC response can be written as
\begin{equation*}
  i \approx g_d v + \frac{g_d}{2 n V_T} v^2 + \cdots,
\end{equation*}
where $g_d$ is the small-signal conductance (the slope of the I--V curve at the bias point).
The first term is linear; the second is quadratic; for small signals, higher-order terms are smaller still.

The quadratic term is the one that matters for both detection and mixing---but for different reasons in each case.

\textbf{Detection (single tone).}
When an RF signal rides on the diode, the quadratic term produces a DC component proportional to the signal amplitude squared---that is, proportional to the RF power.
This is exactly what ``square-law detector'' means: DC output is proportional to RF power input.
It is why a crystal radio converts an AM signal's envelope into audio current.

\textbf{Mixing (two tones).}
When two signals at frequencies $f_1$ and $f_2$ are both present, the quadratic term cross-multiplies them and produces outputs at $f_1 + f_2$ and $f_1 - f_2$.
The linear piece contributes nothing at the sum or difference frequency---it passes both signals through unchanged.

\section*{The catch: the linear piece always wins}

Here is what the textbooks gloss over.
The linear piece---the one that contributes nothing useful for detection or mixing---is always present, and for small signals it is always \emph{larger} than the quadratic piece.

Think of it this way.
Near any operating point, the diode curve is smooth and can be approximated locally as a straight line---a slope, or conductance.
The quadratic term is the \emph{curvature} of that line.
Curvature is always smaller than slope for small signals; it only becomes significant when the signal is large enough that the curve bends noticeably over the full voltage swing.
Looking back at the equation, the quadratic term carries an extra factor of $v / 2nV_T$ relative to the linear term, so any signal much smaller than $2nV_T \approx \SI{50}{\milli\volt}$ is dominated by the linear piece.

So for very small signals the diode is effectively linear.
For detection this means no DC output: the signal passes through the linear conductance and nothing gets rectified.
For mixing it means no frequency conversion: the two tones pass through without producing sum or difference products.

The diode is always more linear than it is square-law.

\section*{The sweet spot}

There is a range of signal amplitudes where the quadratic term is the dominant nonlinearity, even though the linear term is still larger overall.
For a germanium or Schottky diode at room temperature this range spans roughly 1\,mV to 25\,mV across the junction.
For a silicon signal diode it is somewhat wider, but reached at higher absolute signal levels.

Within this range the square-law approximation is genuinely useful:
\begin{itemize}
  \item Detector DC output is proportional to power---clean and predictable.
  \item Mixer conversion is dominated by $f_1 \pm f_2$, with higher-order products small enough to manage.
\end{itemize}

Below this range the diode is effectively linear: detection sensitivity collapses and mixers produce no conversion.
Above it, higher-order terms take over: detector output no longer tracks power cleanly, and mixers generate intermodulation products beyond $f_1 \pm f_2$.

\section*{Practical consequences}

\textbf{Why Ge and Schottky beat Si for weak-signal detection.}
A germanium or zero-bias Schottky diode reaches its square-law range at lower absolute signal levels than a silicon p--n diode, because its forward voltage is lower.
The physics is the same in each case; the threshold is simply reached with less signal power.
This is the fundamental reason crystal radio builders and microwave detector designers reach for Ge and Schottky diodes.

\textbf{Why bias helps.}
A small forward bias current moves the operating point up onto the exponential curve, where both the conductance and the curvature are higher.
This shifts the square-law range to lower signal levels, making a biased detector more sensitive than an unbiased one---at the cost of a bias circuit and quiescent current drain.

\textbf{LO drive level in mixers.}
A mixer driven too hard is no longer operating in the square-law regime.
Higher-order terms produce additional intermodulation products and compress the IF output---the converted signal level stops rising in proportion to the input.
``Square-law mixer'' behavior is only valid over a moderate range of LO drive levels.

\section*{The bottom line}

A diode is a square-law device in the same way a stretch of road is flat: it is a useful approximation over a limited range, and knowing the limits prevents you from driving off the edge.
The exponential is the truth; the square law is a convenience that holds over roughly one decade of signal amplitude.
Understanding that range---and what sits on either side of it---is what separates a detector that works from one that merely seems to.

\end{document}
